\section{Hierarchical Data and the Closure Property}
\label{Section 2.2}

\noindent
Before beginning this section we should say what shape the next few subsections will have,
because this is the place where the two languages start to truly differ.

The pair is \emph{the} fundamental data structure of Lisp.  Everything compound
is built from it: a list is a chain of pairs, a tree is a also chain of pairs
with different nesting, and so is the text of the program itself.  The Python tuple is \emph{a} fundamental data
structure of Python, but not \emph{the} one---and it is not confined to two
elements.  The pair we built in \link{Section 2.1} was simply a tuple that
happened to have two, a special case rather than the foundation.

So where the original constructs a list by chaining pairs together, we shall use
a tuple directly.  The chain is still worth knowing, and we shall meet it in
\link{Section 2.2.1}: it is the right way to think about a list while writing a
recursive procedure, and we shall start by writing our procedures recursively for the
thought process it promotes.  But we won't build on it, because that is not how Python
is written.

The introduction to this section is given in the original's terms, pairs and \code{cons} and
all, and the diagrams are the original's diagrams.  Read \code{(cons a b)} as the
two-element tuple \code{(a, b)}; the argument being made about it holds for
tuples of any length.

As we have seen, pairs provide a primitive ``glue'' that we can use to
construct compound data objects.  \link{Figure 2.2} shows a standard way to
visualize a pair---in this case, the pair formed by \code{(cons 1 2)}.  In this
representation, which is called \newterm{box-and-pointer notation}, each object
is shown as a \newterm{pointer} to a box.  The box for a primitive object
contains a representation of the object.  For example, the box for a number
contains a numeral.  The box for a pair is actually a double box, the left part
containing (a pointer to) the \code{car} of the pair and the right part
containing the \code{cdr}.

We have already seen that \code{cons} can be used to combine not only numbers
but pairs as well.  (You made use of this fact, or should have, in doing
\link{Exercise 2.2} and \link{Exercise 2.3}.)  As a consequence, pairs provide a
universal building block from which we can construct all sorts of data
structures.  \link{Figure 2.3} shows two ways to use pairs to combine the
numbers 1, 2, 3, and 4.


% figure 2.2
\begin{imageFigure}{Box-and-pointer representation of \code{(cons 1 2)}}
  % Box-and-pointer representation of (cons 1 2), for Section 2.2.
% Redrawn from upstream's figure 2.2 so the web edition has a picture at all.
% Shared styles for the box-and-pointer figures of Chapter 2.
%
% \input at the top of each such figure, so that the cells, leaves, pointers
% and nil strokes are drawn identically wherever they appear. A pair is ONE
% rounded rectangle with a divider through it, never two abutting boxes:
% rounding two boxes would round their inner edges too, which is not how the
% original's cells look.
\tikzset{
  pair/.style  = {draw, rounded corners=1.2pt, fill=black!5, minimum width=16mm,
                  minimum height=7mm, inner sep=0pt, outer sep=0pt},
  half/.style  = {minimum width=8mm, minimum height=7mm,
                  inner sep=0pt, outer sep=0pt},
  leaf/.style  = {draw, rounded corners=1.2pt, fill=white, minimum width=8mm,
                  minimum height=7mm, inner sep=0pt, outer sep=0pt,
                  font=\ttfamily\footnotesize},
  dot/.style   = {circle, fill, inner sep=1.1pt},
  ptr/.style   = {-{Latex[length=1.6mm, width=1.4mm]}, thin},
  plabel/.style = {font=\ttfamily\scriptsize, inner sep=1pt},
}

\begin{tikzpicture}[y = -1mm, x = 1mm]

\node[half] (pl) at (8, 0) {};
\node[half] (pr) at (16, 0) {};
\node[pair] (p)  at (12, 0) {};
\draw (pl.north east) -- (pl.south east);
\node[dot] (pld) at (8, 0) {};
\node[dot] (prd) at (16, 0) {};

\node[leaf] (l1) at (4, 14) {1};
\node[leaf] (l2) at (20, 14) {2};

\draw[ptr] (pld) -- (l1.north);
\draw[ptr] (prd) -- (l2.north);

\end{tikzpicture}

\end{imageFigure}

% figure 2.3
\begin{imageFigure}{Two ways to combine 1, 2, 3, and 4 using pairs.}
  \resizebox{\linewidth}{!}{% Two ways to combine 1, 2, 3 and 4 using pairs, for Section 2.2.
% Redrawn from upstream's figure 2.3.
% Shared styles for the box-and-pointer figures of Chapter 2.
%
% \input at the top of each such figure, so that the cells, leaves, pointers
% and nil strokes are drawn identically wherever they appear. A pair is ONE
% rounded rectangle with a divider through it, never two abutting boxes:
% rounding two boxes would round their inner edges too, which is not how the
% original's cells look.
\tikzset{
  pair/.style  = {draw, rounded corners=1.2pt, fill=black!5, minimum width=16mm,
                  minimum height=7mm, inner sep=0pt, outer sep=0pt},
  half/.style  = {minimum width=8mm, minimum height=7mm,
                  inner sep=0pt, outer sep=0pt},
  leaf/.style  = {draw, rounded corners=1.2pt, fill=white, minimum width=8mm,
                  minimum height=7mm, inner sep=0pt, outer sep=0pt,
                  font=\ttfamily\footnotesize},
  dot/.style   = {circle, fill, inner sep=1.1pt},
  ptr/.style   = {-{Latex[length=1.6mm, width=1.4mm]}, thin},
  plabel/.style = {font=\ttfamily\scriptsize, inner sep=1pt},
}

\begin{tikzpicture}[y = -1mm, x = 1mm]

% --- left: (cons (cons 1 2) (cons 3 4)) ---
\foreach \n/\x/\y in {R/10/0, A/6/17, B/40/0} {
  \node[half] (\n l) at (\x, \y) {};
  \node[half] (\n r) at (\x+8, \y) {};
  \node[pair] (\n)   at (\x+4, \y) {};
  \draw (\n l.north east) -- (\n l.south east);
  \node[dot] (\n d)  at (\x, \y) {};
  \node[dot] (\n rd) at (\x+8, \y) {};
}
\node[leaf] (L1) at (2, 34) {1};
\node[leaf] (L2) at (16, 34) {2};
\node[leaf] (L3) at (38, 17) {3};
\node[leaf] (L4) at (52, 17) {4};
\draw[ptr] (Rd)  -- (A.north);
\draw[ptr] (Rrd) -- (Bl.west);
\draw[ptr] (Ad)  -- (L1.north);
\draw[ptr] (Ard) -- (L2.north);
\draw[ptr] (Bd)  -- (L3.north);
\draw[ptr] (Brd) -- (L4.north);
\node[plabel] at (26, 46) {(cons (cons 1 2) (cons 3 4))};

% --- right: (cons (cons 1 (cons 2 3)) 4) ---
\foreach \n/\x/\y in {S/92/0, C/88/17, D/112/17} {
  \node[half] (\n l) at (\x, \y) {};
  \node[half] (\n r) at (\x+8, \y) {};
  \node[pair] (\n)   at (\x+4, \y) {};
  \draw (\n l.north east) -- (\n l.south east);
  \node[dot] (\n d)  at (\x, \y) {};
  \node[dot] (\n rd) at (\x+8, \y) {};
}
\node[leaf] (M4) at (120, 0) {4};
\node[leaf] (M1) at (86, 34) {1};
\node[leaf] (M2) at (110, 34) {2};
\node[leaf] (M3) at (124, 34) {3};
\draw[ptr] (Sd)  -- (C.north);
\draw[ptr] (Srd) -- (M4.west);
\draw[ptr] (Cd)  -- (M1.north);
\draw[ptr] (Crd) -- (Dl.west);
\draw[ptr] (Dd)  -- (M2.north);
\draw[ptr] (Drd) -- (M3.north);
\node[plabel] at (106, 46) {(cons (cons 1 (cons 2 3)) 4)};

\end{tikzpicture}
}
\end{imageFigure}

The ability to create pairs whose elements are pairs is the essence of list
structure's importance as a representational tool.  We refer to this ability as
the \newterm{closure property} of \code{cons}.  In general, an operation for
combining data objects satisfies the closure property if the results of
combining things with that operation can themselves be combined using the same
operation.\footnote{The use of the word ``closure'' here comes from abstract
algebra, where a set of elements is said to be closed under an operation if
applying the operation to elements in the set produces an element that is again
an element of the set.  The Lisp community also (unfortunately) uses the word
``closure'' to describe a totally unrelated concept: a closure is an
implementation technique for representing procedures with free variables.  The
original declines to use the word in that second sense.  We cannot: it is the
ordinary word for the thing in Python, and we have used it so already---in
\link{Section 1.1.8}, where a procedure returned by another kept the bindings it
was made with, and again in \link{Section 2.1.3}, where we looked at
\code{\_\_closure\_\_} and read those bindings off directly.  The two senses really
are unrelated, and the reader must simply carry both.  In this section the word
means the one from algebra.} Closure
is the key to power in any means of combination because it permits us to create
\newterm{hierarchical} structures---structures made up of parts, which
themselves are made up of parts, and so on.

From the outset of \link{Chapter 1}, we've made essential use of closure in
dealing with procedures, because all but the very simplest programs rely on the
fact that the elements of a combination can themselves be combinations.  In
this section, we take up the consequences of closure for compound data.  We
describe some conventional techniques for using pairs to represent sequences
and trees, and we exhibit a graphics language that illustrates closure in a
vivid way.\footnote{The notion that a means of combination should satisfy
closure is a straightforward idea.  Unfortunately, the data combiners provided
in many popular programming languages do not satisfy closure, or make closure
cumbersome to exploit.  In Fortran or Basic, one typically combines data
elements by assembling them into arrays---but one cannot form arrays whose
elements are themselves arrays.  Pascal and C admit structures whose elements
are structures.  However, this requires that the programmer manipulate pointers
explicitly, and adhere to the restriction that each field of a structure can
contain only elements of a prespecified form.  Unlike Lisp with its pairs,
these languages have no built-in general-purpose glue that makes it easy to
manipulate compound data in a uniform way.  Python is not among the accused: a
tuple may contain tuples, and a list lists, without restriction or ceremony, so
Python satisfies closure as squarely as Lisp does.  This limitation lies behind Alan
Perlis's comment in his foreword to this book: ``In Pascal the plethora of
declarable data structures induces a specialization within functions that
inhibits and penalizes casual cooperation.  It is better to have 100 functions
operate on one data structure than to have 10 functions operate on 10 data
structures.''}



